\section*{अध्याय \ref{ch:b3:homogeneous-catalysis} --- उद्योग में समांगी उत्प्रेरण}

\begin{solution}{exo:b3:homogeneous-catalysis:1}
$\ce{RhCl(PPh3)3}$: 16, Rh(I) $d^8$। $\ce{RhCl(PPh3)2}$: 14, Rh(I)। $\ce{RhClH2(PPh3)2}$: 16, Rh(III) $d^6$।
ऐल्कीन संकुल: 18, Rh(III)। ऐल्किल हाइड्राइड: 16, Rh(III)। अपचायी विलोपन
14, Rh(I) पर लौटा देता है।
\end{solution}

\begin{solution}{exo:b3:homogeneous-catalysis:2}
ब्यूटेनैल, $\ce{CH3CH2CH2CHO}$ (रैखिक), और 2-मेथिलप्रोपेनैल, $\ce{(CH3)2CHCHO}$ (शाखित)।
\end{solution}

\begin{solution}{exo:b3:homogeneous-catalysis:3}
डाइएथिल साइक्लोपेन्ट-3-ईन-1,1-डाइकार्बोक्सिलेट, एक पंचसदस्यीय वलय, एथीन के
निकास के साथ: दोनों अंतस्थ $\ce{CH2}$ समूह $\ce{C2H4}$ के रूप में निकल जाते हैं और दोनों आंतरिक CH कार्बन
जुड़ जाते हैं।
\end{solution}

\begin{solution}{exo:b3:homogeneous-catalysis:4}
मेथिल (\textit{E})-3-फ़ेनिलप्रोप-2-ईनोएट (मेथिल सिनामेट): ऐरिल
अंतस्थ कार्बन से जुड़ता है, और अधिक स्थायी संरूपी से सिन $\beta$-हाइड्राइड विलोपन
\textit{E} ऐल्कीन देता है।
\end{solution}

\begin{solution}{exo:b3:homogeneous-catalysis:5}
ऑक्सीकारी योग धीमा है और $\ce{[Rh(CO)2I2]-}$ विराम अवस्था है, अतः
$v = k[\mathrm{Rh}]_{\mathrm{tot}}[\ce{CH3I}]$। मेथेनॉल केवल HI के साथ एक तीव्र अभिक्रिया द्वारा मेथिल आयोडाइड को पुनर्जनित करता है;
डाला गया कुल आयोडाइड $[\ce{CH3I}]$ तय करता है, अतः दर मेथेनॉल
सांद्रता पर तब तक निर्भर नहीं करती जब तक HI को वापस मेथिल आयोडाइड में बदलने के लिए
बहुत कम मेथेनॉल न बचे।
\end{solution}

\begin{solution}{exo:b3:homogeneous-catalysis:6}
4-ब्रोमोटॉलूईन के साथ फ़ेनिलबोरोनिक अम्ल, या ब्रोमोबेंज़ीन के साथ
4-मेथिलफ़ेनिलबोरोनिक अम्ल। दोनों काम करते हैं; पहला सबसे सस्ता बोरोनिक अम्ल और एक स्थायी,
आम ऐरिल ब्रोमाइड उपयोग करता है।
\end{solution}

\begin{solution}{exo:b3:homogeneous-catalysis:7}
$C_2$-सममित: आइसोटैक्टिक (दोनों स्थल एक ही फलक चुनते हैं)। $C_s$-सममित:
सिंडियोटैक्टिक (दर्पण-प्रतिबिंब स्थल एकांतर होते हैं)। असेतुबद्ध $\ce{Cp2ZrCl2}$: ऐटैक्टिक (अकिरैल
स्थल)।
\end{solution}

\begin{solution}{exo:b3:homogeneous-catalysis:8}
$\mathrm{TON} = 0.98/0.0010 = 980$; माध्य $\mathrm{TOF} = 980/\qty{2.0}{h} = \qty{490}{h^{-1}}$।
\end{solution}

\begin{solution}{exo:b3:homogeneous-catalysis:9}
पुनरावृत्त इकाई $\ce{-[CH=CH-C5H8]-}$ है, जिसमें दोनों CH समूह एक साइक्लोपेन्टेन वलय के कार्बन 1 और 3
पर हैं (पॉली(1,3-साइक्लोपेन्टिलीनवाइनिलीन))। द्विचक्रीय वलय के
तनाव का मोचन $\Delta_rH^\circ$ को प्रबल रूप से ऋणात्मक बनाता है, जो स्थानांतरीय एन्ट्रॉपी की
हानि से भारी पड़ता है; कोई गैस नहीं बनती, अतः प्रेरक बल
एन्थैल्पीय है।
\end{solution}

\begin{solution}{exo:b3:homogeneous-catalysis:10}
$\bar n_n = 1/(1 - p) = 2000$; $M_w/M_n = 1 + p = 1.9995 \approx 2.0$। \omterm{def:b3:homogeneous-catalysis:ziegler-natta}{ज़ीग्लर--नाटा उत्प्रेरक} में
अनेक प्रकार के स्थल होते हैं, प्रत्येक का अपना $p$; भिन्न माध्यों वाले अनेक शुल्ज़--फ़्लोरी
वितरणों का योग अधिक चौड़ा होता है (परिक्षेपिता प्रायः 4 से 8)।
\end{solution}

\begin{solution}{exo:b3:homogeneous-catalysis:11}
निम्न दाब पर $K_1K_2[\ce{H2}] \ll [\mathrm L] + K_1$: $v \approx kK_1K_2[\mathrm{Rh}]_{\mathrm{tot}}[\text{ऐल्कीन}][\ce{H2}]/([\mathrm L] + K_1)$,
$\ce{H2}$ में प्रथम कोटि। उच्च दाब पर पद $K_1K_2[\ce{H2}]$ प्रभावी होता है: $v \to k[\mathrm{Rh}]_{\mathrm{tot}}[\text{ऐल्कीन}]$,
$\ce{H2}$ से स्वतंत्र। $[\mathrm L]$ केवल हर में आता है: मिलाया गया फ़ॉस्फ़ीन
रोडियम को वापस संतृप्त पूर्व-उत्प्रेरक में धकेलकर $v$ घटा देता है।
\end{solution}

\begin{solution}{exo:b3:homogeneous-catalysis:12}
जो निवेशन रोडियम को आंतरिक कार्बन पर रखता है, वह भीड़ भरी धातु के बगल में एक द्वितीयक ऐल्किल
बनाता है; भारी फ़ॉस्फ़ीन (और उनका आधिक्य, जो धातु पर दो को
बनाए रखता है) उस संक्रमण अवस्था को प्राथमिक ऐल्किल देने वाली संक्रमण अवस्था से
बदतर बना देते हैं, जो रैखिक ऐल्डिहाइड तक ले जाती है। रैखिक ब्यूटेनैल ही वह है जिसे
ऐल्डोल संघनन और हाइड्रोजनन द्वारा सुघट्यकारियों में प्रयुक्त शाखित $C_8$ ऐल्कोहॉल में बदला जाता है;
2-मेथिलप्रोपेनैल एक कम-मूल्य उप-उत्पाद है।
\end{solution}

\begin{solution}{pb:b3:homogeneous-catalysis:1}
\textbf{1.} Ir(I), $d^8$: $9 + 4 + 2 + 1 = 16$ इलेक्ट्रॉन (उदासीन गणना, आवेश सहित)।

\textbf{2.} $\ce{[CH3Ir(CO)2I3]-}$, 18 इलेक्ट्रॉन, Ir(III)।

\textbf{3.} $\ce{[CH3Ir(CO)3I2]}$, 18 इलेक्ट्रॉन, Ir(III)।

\textbf{4.} $\ce{[CH3COIr(CO)2I2]}$, 16 इलेक्ट्रॉन, Ir(III)।

\textbf{5.} $\ce{[CH3COIr(CO)2I3]-}$ (18) $\ce{CH3COI}$ का विलोपन करता है, जिससे $\ce{[Ir(CO)2I2]-}$ पुनर्जनित होता है।

\textbf{6.} $\ce{CH3COI + H2O -> CH3COOH + HI}$; $\ce{CH3OH + HI -> CH3I + H2O}$।

\textbf{7.} $\ce{CH3OH + CO -> CH3COOH}$।

\textbf{8.} इरिडियम के साथ प्रवासी निवेशन धीमा है; मेथिल आयोडाइड का ऑक्सीकारी योग,
जो रोडियम के साथ दर-निर्धारक है, अधिक
इलेक्ट्रॉन-समृद्ध इरिडियम पर कहीं तेज़ है।

\textbf{9.} वे $\ce{[CH3Ir(CO)2I3]-}$ से एक आयोडाइड हटाते हैं, जिससे उदासीन ट्राइकार्बोनिल बनता है, जिसमें
धातु कम पश्च-दान करती है और मेथिल अधिक इलेक्ट्रॉनरागी CO पर
तेज़ी से प्रवास करता है।

\textbf{10.} मुक्त आयोडाइड सांद्रता पर व्युत्क्रम निर्भरता: आयोडाइड
पूर्व-साम्य को वापस धीमे ऋणायनी संकुल की ओर धकेलता है।

\textbf{11.} मेथेनॉल केवल दर-निर्धारक चरण के बाद प्रवेश करता है, HI के साथ एक तीव्र अभिक्रिया द्वारा
मेथिल आयोडाइड में बदलकर।

\textbf{12.} $M(\ce{CH3COOH}) = \qty{60.1}{g/mol}$: $n = \qty{5.0e11}{g}/\qty{60.1}{g/mol} = \qty{8.3e9}{mol}$ प्रति वर्ष।

\textbf{13.} $\qty{8.32e9}{mol}/\qty{8000}{h} = \qty{1.04e6}{mol/h}$।

\textbf{14.} $\qty{2.0e5}{L} \times \qty{5.0e-3}{mol/L} = \qty{1.0e3}{mol}$ इरिडियम, $\qty{1.0e3}{mol} \times \qty{192.2}{g/mol} \approx \qty{190}{kg}$।

\textbf{15.} $\mathrm{TOF} = \qty{1.04e6}{mol/h}/\qty{1.0e3}{mol} = \qty{1.0e3}{h^{-1}}$, लगभग \qty{0.29}{s^{-1}}।

\textbf{16.} $\mathrm{TON} = \qty{8.3e9}{mol}/\qty{1.0e3}{mol} = \num{8.3e6}$ प्रति वर्ष।

\textbf{17.} $\qty{5.0e8}{kg}/(\qty{8000}{h} \times \qty{200}{m^3}) \approx \qty{310}{kg.m^{-3}.h^{-1}}$।

\textbf{18.} $\ce{CO + H2O -> CO2 + H2}$।

\textbf{19.} प्रति मोल ऐसीटिक अम्ल $1/0.94 = \qty{1.064}{mol}$ CO।

\textbf{20.} प्रति मोल ऐसीटिक अम्ल $1.064 - 1 = \qty{0.064}{mol}$ \ce{CO2}।

\textbf{21.} एक टन $\qty{1.0e6}{g}/\qty{60.1}{g/mol} = \qty{1.66e4}{mol}$ है; $\qty{1.66e4}{mol} \times 0.0638 \times \qty{44.0}{g/mol} \approx \qty{47}{kg}$ \ce{CO2}।

\textbf{22.} जल विस्थापन का एक अभिकारक है: कम जल इसे धीमा करता है और अम्ल को सुखाने की
ऊर्जा बचाता है। ऐसीटिल आयोडाइड के जल-अपघटन के लिए और
उत्प्रेरक को सक्रिय और घुला हुआ रखने के लिए जल अब भी आवश्यक है; रोडियम के साथ कम जल रोडियम
आयोडाइड को अवक्षेपित कर देता है।

\textbf{23.} हाइड्रोजन: यह गैस में संचित होता है, CO का आंशिक दाब घटाता है
(अतः गैस को बाहर निकालना पड़ता है, जिससे CO की हानि होती है), और मध्यवर्तियों का हाइड्रोजनन करके
प्रोपियोनिक अम्ल जैसे उप-उत्पाद बना सकता है।

\textbf{24.} इरिडियम उत्प्रेरक लगभग $\qty{1.0e3}{h^{-1}}$ की दर से आवर्त पूरे करता है: प्रत्येक इरिडियम परमाणु
प्रति घंटे ऐसीटिक अम्ल के एक हज़ार अणु बनाता है, प्रति वर्ष लगभग अस्सी लाख।
\end{solution}
